\scanpage{61}\mspage{60}
Or ça n'a pas l'air d'impliquer la condition.
Car ça n'implique pas que \deleted{la disjonction}
la disjonction ensembliste implique la disjonction géométrique. Prenons p.ex.
$M$ réduit à \deleted{\unsure{trois}} \inserted{des} multistrates \deleted{\unsure{chacune}}
réduites à une seule \deleted{ensemble} \inserted{strate} (ce qui
rend vide la première condition \unsure{Atens tf|Atens t.f|Atensfig} 3)
par exemple définies par les \deleted{\unsure{d|t}}trois ensembles
$A,B$ et $\unsure{A'|A}$, $A'\subsetneq A,\ A\cap B=\varnothing$. \deleted{Prenons}
\deleted{\unsure{donc}} $X=\{A\},\quad Y=\{B\},\quad X'=\{A'\}$, et soit $Y'=Y$,
donc $X'\deleted{\unsure{<|\leq}}\ll X,\quad Y'=Y$ donc $Y'\ll Y$. \deleted{On définit} \inserted{La}
\deleted{$\underset{M}{\mathrel{\tikz[baseline=-.45ex,x=1em,y=1em]{\draw[line width=.055em] (.52,.27)--(.05,.055)--(.48,-.11); \draw[line width=.055em] (.14,.21)--(.46,.04)--(0,-.17);}}}$ \unsure{comme étant la|comme étant la co}} \inserted{compatibilité ensembliste} des trois
multistrates est déjà acquise. Donc on
peut prendre pour $\underset{M}{\mathrel{\tikz[baseline=-.45ex,x=1em,y=1em]{\draw[line width=.055em] (.52,.27)--(.05,.055)--(.48,-.11); \draw[line width=.055em] (.14,.21)--(.46,.04)--(0,-.17);}}}$ n'importe quelle
relation \unsure{sym. refl|sym. réfl|sym. rfl}, \deleted{\unsure{donc}} \inserted{correspondant à} n'importe quelle
partie \deleted{\unsure{à deux éléments}} \deleted{\unsure{des}}\unsure{de} l'ens des
trois paires de \deleted{strates} multistrates $X,Y,X'$
on choisira $X\underset{M}{\mathrel{\tikz[baseline=-.45ex,x=1em,y=1em]{\draw[line width=.055em] (.52,.27)--(.05,.055)--(.48,-.11); \draw[line width=.055em] (.14,.21)--(.46,.04)--(0,-.17);}}}Y$, et $X'\underset{M}{\overline{\mathrel{\tikz[baseline=-.45ex,x=1em,y=1em]{\draw[line width=.055em] (.52,.27)--(.05,.055)--(.48,-.11); \draw[line width=.055em] (.14,.21)--(.46,.04)--(0,-.17);}}}}Y$ --
alors $X\underset{\unsure{\mathcal F|F}}{\parallel}Y$, mais $X'\underset{\mathcal F}{\overline{\parallel}}Y$, ok.

On posera donc
\begingroup\makeatletter\tagsleft@true\makeatother
\[
\begin{aligned}
&\text{\deleted{\uline{At} \uline{9}} Si }\ X,Y\ \deleted{\unsure{\in M}}\,,\ X',Y'\in M,\\
&\text{avec }\ X'\ll X,\quad Y'\ll Y,\quad \text{et si }X\parallel Y,\\
&\text{alors }\ X'\parallel Y'\quad(\text{il suffit de poser }X \mathrel{\tikz[baseline=-.45ex,x=1em,y=1em]{\draw[line width=.055em] (.52,.27)--(.05,.055)--(.48,-.11); \draw[line width=.055em] (.14,.21)--(.46,.04)--(0,-.17);}} Y)
\end{aligned}
\tag*{($*$)}
\]
\endgroup
